\fsection{Marco}

El proyecto se enmarca en una iniciativa de innovación e investigación aplicada promovida por el Ministerio
de Educación y Formación Profesional, con financiación de la Unión Europea a través de los fondos Next
Generation EU. En él participan tres centros educativos de Formación Profesional: LH FP Barakaldo, el IES
Javier García Téllez de Cáceres y el IES Jorge Manrique de Motilla del Palancar, junto con la empresa
Fluitronic S.L., que aportó la formación del profesorado y el suministro del equipamiento. Cada centro ha
desarrollado su propio prototipo de forma independiente, compartiendo conocimiento y experiencias a lo
largo del proceso.
